\subsubsection{Fourier Series}\label{subsec:fourier-series-hilbert-bases}

The most important example of expansions with respect to an orthonormal basis is the Fourier expansion.

\begin{Theorem}
\label{thm:fourier-cosine-hilbert-basis}
Let $\phi_0(x)=1$, $\phi_n(x)=\sqrt2\cos(n\pi x)$, $n\in\mathbb{N}$. Then the functions $\phi_n$, $n=0,1,2,\ldots$, form a basis of $L^2(0,1)$.
\end{Theorem}

\begin{Proof}
An easy calculation shows that the $\phi_n$ are an orthonormal system. By Theorem~\ref{thm:hilbert-basis-equivalent-characterizations} it therefore suffices to show that the linear span of the $\phi_n$ is dense in $L^2(0,1)$. Since $C([0,1])$ is dense in $L^2(0,1)$, we only need to show that every continuous function can be approximated by a linear combination of the $\phi_n$. We make the substitution $\cos\pi x=u$, which bijectively maps $[0,1]$ to $[-1,1]$. By the Weierstrass approximation theorem, every continuous function on $[-1,1]$ can be approximated uniformly by polynomials; hence every continuous function on $[0,1]$ can be approximated uniformly by polynomials in $\cos\pi x$. Elementary trigonometric identities show that any expression $\sum_{n=0}^N a_n(\cos\pi x)^n$ can be rewritten in the form $\sum_{n=0}^N b_n\cos(n\pi x)$.
\end{Proof}

Functions in $L^2(0,1)$ can also be expanded in terms of a sine series instead of a cosine series.

\begin{Theorem}
\label{thm:fourier-sine-hilbert-basis}
Let $\psi_n(x)=\sqrt2\sin(n\pi x)$, $n\in\mathbb{N}$. Then the functions $\psi_n$, $n\in\mathbb{N}$, form a basis of $L^2(0,1)$.
\end{Theorem}

\begin{Proof}
We use the fact that $\mathcal{D}(0,1)$ is dense in $L^2(0,1)$. If $f\in\mathcal{D}(0,1)$, then $f(x)/\sin(\pi x)$ is continuous on $[0,1]$ and from the proof of the last theorem we conclude that it can be uniformly approximated by expressions of the form $\sum_{n=0}^N a_n\cos(n\pi x)$. Hence $f(x)$ can be uniformly approximated by expressions of the form
\begin{equation}\label{eq:fourier-sine-cosine-product-identity}
\begin{aligned}
\sum_{n=0}^N a_n\cos(n\pi x)\sin(\pi x)
&=\dfrac12\sum_{n=0}^N a_n\bigl(\sin((n+1)\pi x)\\
&\hspace{8em}-\sin((n-1)\pi x)\bigr).
\end{aligned}
\end{equation}
This completes the proof.
\end{Proof}

Theorems~\ref{thm:fourier-cosine-hilbert-basis} and~\ref{thm:fourier-sine-hilbert-basis} yield the following simple consequence.

\begin{Corollary}
\label{cor:fourier-complex-hilbert-basis}
The functions $(1/\sqrt2)e^{in\pi x}$, $n\in\mathbb{Z}$, form a basis of $L^2(-1,1)$.
\end{Corollary}

\begin{Proof}
Any function in $L^2(-1,1)$ can be decomposed into an even and an odd part. Using the preceding two theorems, we can expand the even part in a cosine series and the odd part in a sine series.
\end{Proof}
